\section{The source calculus \texorpdfstring{$\CICm$}{CIC-}}
\label{sec:calculus}

This section defines the source calculus of the translation: a monomorphic,
first-order fragment of the Calculus of Inductive Constructions.  The
calculus is presented in a \emph{stratified} style: instead of a single
syntactic category of terms classified by sorts, we keep separate
categories for set types, propositions, predicates, terms and proofs.  The
stratification builds the outcome of CIC's sort analysis --- which subterms
are computational, which are logical, which are type-level --- into the
grammar, so that the erasure function of \cref{sec:lambdabox} and the
translation of \cref{sec:translation} are definable by structural recursion
on raw syntax, with no typing oracle.  \Cref{sec:calculus-cic} discusses how
the fragment is carved out of full CIC.

\subsection{Syntax}
\label{sec:calculus-syntax}

We fix countably infinite, pairwise disjoint sets of \emph{term
variables} (written $x, y, z, w, f, \dots$), of \emph{proof variables}
($p$, $h, \dots$) and of \emph{predicate variables} ($q, \dots$), and
disjoint sets of names for \emph{datatypes} ($I$), for \emph{datatype
constructors} ($\wcns{c}$, $\wcns{d}$), for \emph{inductive predicates}
($J$), for their \emph{clause constructors} ($\wcns{k}$), and for
\emph{constants} ($c$).

\begin{definition}[Syntax of $\CICm$]
\label{def:syntax}
The syntactic categories of \emph{set types} $\sigma, \tau$,
\emph{propositions} $\varphi, \psi$, \emph{predicate kinds} $K$,
\emph{predicate expressions} $P$, \emph{terms} $t, a, b$ and \emph{proofs}
$\pi$ are defined by the following grammar.
\begin{align*}
  \sigma, \tau &\Coloneqq
    I
    \mid \Pi x{:}\sigma.\,\tau
    \mid \varphi \to \tau
    \mid \refty{x{:}\sigma}{\varphi}
  \\
  \varphi, \psi &\Coloneqq
    \top \mid \bot
    \mid \varphi \wedge \psi
    \mid \varphi \vee \psi
    \mid \varphi \to \psi
    \mid \forall x{:}\sigma.\,\varphi
    \mid \exists x{:}\sigma.\,\varphi
  \\ &\mathrel{\phantom{\Coloneqq}}
    \mid \forall q{:}K.\,\varphi
    \mid a =_{\sigma} b
    \mid J(\tup{a})
    \mid q(\tup{a})
  \\
  K &\Coloneqq (\sigma_1, \dots, \sigma_n) \to \Prop \qquad (n \geq 1)
  \\
  P &\Coloneqq q \mid \lambda(x_1{:}\sigma_1, \dots, x_n{:}\sigma_n).\,\varphi
  \\
  t, a, b &\Coloneqq
    x
    \mid c
    \mid t\,a
    \mid t\,\pi
    \mid \lambda x{:}\sigma.\,t
    \mid \lambda p{:}\varphi.\,t
    \mid \wcns{c}(\tup{a})
  \\ &\mathrel{\phantom{\Coloneqq}}
    \mid \tcase{a}{z.\tau}{\tup{br}}
    \mid \tfix\,f{:}T.\,t
    \mid \trefine{t}{\pi}
    \mid \tval{t}
  \\ &\mathrel{\phantom{\Coloneqq}}
    \mid \tsplit{\pi}{p_1\,p_2.\,t}
    \mid \tcast{\pi}{z.\tau}{t}
    \mid \texfalso{\pi}{\sigma}
  \\
  \pi &\Coloneqq
    p
    \mid c
    \mid \lambda p{:}\varphi.\,\pi
    \mid \pi\,\pi'
    \mid \lambda x{:}\sigma.\,\pi
    \mid \pi\,a
    \mid \Lambda q{:}K.\,\pi
    \mid \pi\,P
    \mid \langle \pi_1, \pi_2 \rangle
    \mid \pi.1
    \mid \pi.2
  \\ &\mathrel{\phantom{\Coloneqq}}
    \mid \iota_1(\pi) \mid \iota_2(\pi)
    \mid \mathsf{vcase}(\pi;\, p_1.\pi_1;\, p_2.\pi_2)
    \mid \langle a, \pi \rangle_{\exists x{:}\sigma.\varphi}
    \mid \mathsf{ecase}(\pi;\, x\,p.\,\pi')
  \\ &\mathrel{\phantom{\Coloneqq}}
    \mid \mathsf{tt}
    \mid \texfalso{\pi}{\varphi}
    \mid \mathsf{refl}_{\sigma}(a)
    \mid \mathsf{rew}(\pi;\, z.\psi;\, \pi')
    \mid \wcns{k}(\tup{a}; \tup{\pi})
    \mid \tind{J}(\tup{z}.\psi;\, \tup{\pi};\, \pi)
  \\ &\mathrel{\phantom{\Coloneqq}}
    \mid \tind{I}(z.\psi;\, \tup{\pi};\, a)
    \mid \tvalid{t}
  \end{align*}
where a branch list $\tup{br}$ has the form $\{\wcns{c}_1(\tup{y}_1)
\Rightarrow t_1 \mid \dots \mid \wcns{c}_m(\tup{y}_m) \Rightarrow t_m\}$
with one branch for every constructor of the datatype concerned, and $T$ in
$\tfix\,f{:}T.\,t$ ranges over set types.  In
$\tcase{a}{z.\tau}{\tup{br}}$ the variable $z$ binds into the motive
$\tau$; in $\refty{x{:}\sigma}{\varphi}$ the variable $x$ binds into
$\varphi$; the binding structure of the remaining forms is the evident one.
Terms of the form $\trefine{t}{\pi}$ carry a refinement-type annotation
$\refty{x{:}\sigma}{\varphi}$, and $\langle a, \pi\rangle$ carries its
$\exists$-type; we omit the annotations when they are irrelevant.  We write
$\neg\varphi$ for $\varphi \to \bot$ and
$\varphi \leftrightarrow \psi$ for
$(\varphi \to \psi) \wedge (\psi \to \varphi)$.
Terms, proofs, types, propositions and predicate expressions are identified
up to $\alpha$-equivalence, and substitution is capture-avoiding.
\end{definition}

Three points deserve comment.

First, the calculus has \emph{three} function spaces at the term level:
$\Pi x{:}\sigma.\tau$ (informative argument), $\varphi \to \tau$
(propositional premise; note that the premise does not bind a variable
into $\tau$, so the type is non-dependent, although the term-level
abstraction $\lambda p{:}\varphi.t$ does bind $p$ in the body), and the
proof-level implications and quantifications.  Dependency of types on
terms enters exclusively through refinement types
$\refty{x{:}\sigma}{\varphi}$ and through the propositions embedded in
types --- there are no type-level functions, no universes and no large
elimination, so set types are \emph{rigid}: their outermost former is
stable under conversion (\cref{lem:rigid}).

Second, the singleton eliminations of CIC appear as the primitive term
formers $\tsplit{\pi}{p_1 p_2.t}$ (elimination of a conjunction proof into
a term), $\tcast{\pi}{z.\tau}{t}$ (transport along an equality proof) and
$\texfalso{\pi}{\sigma}$ ($\bot$-elimination into a term).  These are
exactly the propositional eliminations into informative positions that
CIC's elimination restriction permits: at most one constructor, all
arguments propositional~\cite{Rocq,Letouzey2003}.

Third, refinement types come with the constructor $\trefine{t}{\pi}$
(pairing a value with a proof of the payload), the projection $\tval{t}$
(the value) and the proof former $\tvalid{t}$ (the payload); this is
$\mathsf{sig}$/$\mathsf{exist}$/$\mathsf{proj1\_sig}$/%
$\mathsf{proj2\_sig}$ after Letouzey's informative-singleton unboxing
\cite[\S 4.3]{Letouzey2004} has been applied at the level of the calculus
itself.

\subsection{Environments}
\label{sec:calculus-env}

\begin{definition}[Declarations and environments]
\label{def:env}
A \emph{global environment} $\Env$ is a finite list of declarations of the
following three kinds.
\begin{enumerate}[leftmargin=2em]
\item A \emph{datatype declaration}
  $\mathsf{data}\ I \defeq \{\, \wcns{c}_1 : (\tup{\sigma}_1);\ \dots;\
  \wcns{c}_m : (\tup{\sigma}_m) \,\}$ with $m \geq 1$ constructors.  Each
  constructor argument type $\sigma_{ij}$ is either the datatype name $I$
  itself (a \emph{recursive argument}) or a closed set type over the
  preceding environment in which $I$ does not occur.
\item An \emph{inductive-predicate declaration}
  $\mathsf{pred}\ J : (\tup{\sigma}) \to \Prop$ with a finite list of
  clauses
  \[
    \wcns{k}_j : \forall \tup{x}_j{:}\tup{\tau}_j.\;
      \varphi_{j1} \to \dots \to \varphi_{j\ell_j} \to J(\tup{u}_j)
  \]
  where the index types $\tup{\sigma}$ are closed set types over the
  preceding environment, the $\tup{\tau}_j$ are set types over the
  preceding environment (possibly mentioning $\tup{x}_j$ bound earlier in
  the same telescope, through refinements), the indices $\tup{u}_j$ are
  terms over $\tup{x}_j$, and each premise $\varphi_{ji}$ is either of the
  form $J(\tup{v})$ (a \emph{recursive premise}) or a proposition over the
  preceding environment and $\tup{x}_j$ in which $J$ does not occur.
\item A \emph{definition}, either a \emph{term definition}
  $c \defeq t : \sigma$ with $\sigma$ a closed set type and $t$ a closed
  term, or a \emph{lemma} $c \defeq \pi : \varphi$ with $\varphi$ a closed
  proposition and $\pi$ a closed proof.
\end{enumerate}
All names declared in $\Env$ are distinct.  An environment consisting only
of declarations of kinds (1)--(3) --- i.e., in which every constant has a
body --- is called \emph{definitional}; this paper considers definitional
environments only (see \cref{rem:axioms} for postulated axioms).
\end{definition}

\begin{definition}[Contexts]
A \emph{context} $\Gamma$ is a finite list of declarations $x : \sigma$,
$p : \varphi$ and $q : K$ with pairwise distinct variables, each type,
proposition or kind well formed over the preceding declarations
(\cref{fig:wf}).
\end{definition}

\subsection{Typing}
\label{sec:calculus-typing}

The judgements of the calculus are: $\vdash \Env \wf$ (environment
well-formedness); and, relative to a fixed well-formed $\Env$ which we
leave implicit, $\Gamma \vdash \sigma \styp$, $\Gamma \vdash \varphi
\sprop$, $\Gamma \vdash P : K$, $\Gamma \vdash t : \sigma$, and $\Gamma
\vdash \pi : \varphi$.  The rules are given in
\cref{fig:wf,fig:terms,fig:proofs}.  Environment well-formedness requires
each declaration to be well formed over its prefix: the type and
proposition components of datatype and predicate declarations must be well
formed in the respective telescopes, and a definition $c \defeq t : \sigma$
(resp.\ $c \defeq \pi : \varphi$) requires $\vdash \sigma \styp$ and
$\vdash t : \sigma$ (resp.\ $\vdash \varphi \sprop$ and $\vdash \pi :
\varphi$) over the prefix.

\begin{figure}[t]
\small
\[
\begin{array}{c}
\rul{I \in \Env}{\Gamma \vdash I \styp}
\qquad
\rul{\Gamma \vdash \sigma \styp \quad \Gamma, x{:}\sigma \vdash \tau \styp}
    {\Gamma \vdash \Pi x{:}\sigma.\tau \styp}
\qquad
\rul{\Gamma \vdash \varphi \sprop \quad \Gamma \vdash \tau \styp}
    {\Gamma \vdash \varphi \to \tau \styp}
\\[2.2ex]
\rul{\Gamma \vdash \sigma \styp \quad \Gamma, x{:}\sigma \vdash \varphi \sprop}
    {\Gamma \vdash \refty{x{:}\sigma}{\varphi} \styp}
\qquad
\rul{}{\Gamma \vdash \top \sprop}
\qquad
\rul{}{\Gamma \vdash \bot \sprop}
\\[2.2ex]
\rul{\Gamma \vdash \varphi \sprop \quad \Gamma \vdash \psi \sprop}
    {\Gamma \vdash \varphi \mathbin{\circledast} \psi \sprop}
\;{\scriptstyle \circledast \in \{\wedge,\vee,\to\}}
\qquad
\rul{\Gamma \vdash \sigma \styp \quad \Gamma, x{:}\sigma \vdash \varphi \sprop}
    {\Gamma \vdash Qx{:}\sigma.\varphi \sprop}
\;{\scriptstyle Q \in \{\forall,\exists\}}
\\[2.2ex]
\rul{\Gamma, q{:}K \vdash \varphi \sprop \quad \Gamma \vdash K\ \mathsf{kind}}
    {\Gamma \vdash \forall q{:}K.\varphi \sprop}
\qquad
\rul{\Gamma \vdash a : \sigma \quad \Gamma \vdash b : \sigma}
    {\Gamma \vdash a =_{\sigma} b \sprop}
\\[2.2ex]
\rul{J : (\tup{\sigma}) \to \Prop \in \Env \quad \Gamma \vdash \tup{a} : \tup{\sigma}}
    {\Gamma \vdash J(\tup{a}) \sprop}
\qquad
\rul{q : (\tup{\sigma}) \to \Prop \in \Gamma \quad \Gamma \vdash \tup{a} : \tup{\sigma}}
    {\Gamma \vdash q(\tup{a}) \sprop}
\\[2.2ex]
\rul{\Gamma \vdash \sigma_i \styp \ (1 \le i \le n)}
    {\Gamma \vdash (\sigma_1,\dots,\sigma_n) \to \Prop\ \mathsf{kind}}
\qquad
\rul{q : K \in \Gamma}{\Gamma \vdash q : K}
\\[2.2ex]
\rul{\Gamma, x_1{:}\sigma_1, \dots, x_n{:}\sigma_n \vdash \varphi \sprop}
    {\Gamma \vdash \lambda(\tup{x}{:}\tup{\sigma}).\varphi : (\tup{\sigma}) \to \Prop}
\end{array}
\]
\caption{Well-formedness of set types, propositions, kinds and predicate
  expressions.  In the rule for $J(\tup{a})$ (and analogously for
  $q(\tup{a})$ and index telescopes), $\Gamma \vdash \tup{a} :
  \tup{\sigma}$ abbreviates $\Gamma \vdash a_i : \sigma_i$ for all $i$;
  since index types are closed, no substitution occurs in the telescope.}
\label{fig:wf}
\end{figure}

\begin{figure}[t]
\small
\[
\begin{array}{c}
\rul{x : \sigma \in \Gamma}{\Gamma \vdash x : \sigma}
\qquad
\rul{(c \defeq t : \sigma) \in \Env}{\Gamma \vdash c : \sigma}
\qquad
\rul{\Gamma \vdash t : \Pi x{:}\sigma.\tau \quad \Gamma \vdash a : \sigma}
    {\Gamma \vdash t\,a : \subst{\tau}{a}{x}}
\\[2.2ex]
\rul{\Gamma \vdash t : \varphi \to \tau \quad \Gamma \vdash \pi : \varphi}
    {\Gamma \vdash t\,\pi : \tau}
\qquad
\rul{\Gamma, x{:}\sigma \vdash t : \tau}
    {\Gamma \vdash \lambda x{:}\sigma.t : \Pi x{:}\sigma.\tau}
\qquad
\rul{\Gamma, p{:}\varphi \vdash t : \tau}
    {\Gamma \vdash \lambda p{:}\varphi.t : \varphi \to \tau}
\\[2.2ex]
\rul{\wcns{c} : (\sigma_1,\dots,\sigma_r) \in I \quad
     \Gamma \vdash a_i : \sigma_i \ (1 \le i \le r)}
    {\Gamma \vdash \wcns{c}(\tup{a}) : I}
\qquad
\rul{\Gamma \vdash t : \refty{x{:}\sigma}{\varphi}}
    {\Gamma \vdash \tval{t} : \sigma}
\\[2.2ex]
\rul{\Gamma \vdash a : I \quad
     \Gamma, z{:}I \vdash \tau \styp \quad
     \Gamma, \tup{y}_i{:}\tup{\sigma}_i \vdash t_i :
       \subst{\tau}{\wcns{c}_i(\tup{y}_i)}{z} \ \ (\text{all } i)}
    {\Gamma \vdash \tcase{a}{z.\tau}{\{\wcns{c}_i(\tup{y}_i) \Rightarrow
      t_i\}_i} : \subst{\tau}{a}{z}}
\\[2.2ex]
\rul{\Gamma \vdash \refty{x{:}\sigma}{\varphi} \styp \quad
     \Gamma \vdash t : \sigma \quad
     \Gamma \vdash \pi : \subst{\varphi}{t}{x}}
    {\Gamma \vdash \trefine{t}{\pi} : \refty{x{:}\sigma}{\varphi}}
\qquad
\rul{\Gamma \vdash \pi : \bot \quad \Gamma \vdash \sigma \styp}
    {\Gamma \vdash \texfalso{\pi}{\sigma} : \sigma}
\\[2.2ex]
\rul{\Gamma \vdash \pi : \varphi_1 \wedge \varphi_2 \quad
     \Gamma, p_1{:}\varphi_1, p_2{:}\varphi_2 \vdash t : \tau}
    {\Gamma \vdash \tsplit{\pi}{p_1\,p_2.t} : \tau}
\qquad
\rul{\Gamma \vdash t : \sigma \quad \sigma \conv \sigma' \quad
     \Gamma \vdash \sigma' \styp}
    {\Gamma \vdash t : \sigma'}
\\[2.2ex]
\rul{\Gamma \vdash \pi : a =_{\sigma} b \quad
     \Gamma, z{:}\sigma \vdash \tau \styp \quad
     \Gamma \vdash t : \subst{\tau}{a}{z}}
    {\Gamma \vdash \tcast{\pi}{z.\tau}{t} : \subst{\tau}{b}{z}}
\\[2.2ex]
\rul{\begin{array}{c}
     T = \Pi x_1{:}\sigma_1 \dots \Pi x_n{:}\sigma_n.\tau \quad
     \sigma_k = I \quad \Gamma \vdash T \styp \\
     t = \lambda \tup{x}{:}\tup{\sigma}.\,
         \tcase{x_k}{z.\subst{\tau}{z}{x_k}}{\{\wcns{c}_i(\tup{y}_i)
         \Rightarrow b_i\}_i} \\
     \Gamma, f{:}T, \tup{x}{:}\tup{\sigma}, \tup{y}_i{:}\tup{\sigma}'_i
       \vdash b_i : \subst{\tau}{\wcns{c}_i(\tup{y}_i)}{x_k}
       \ \ (\text{all } i) \quad
     \mathrm{grd}_k(f; \tup{y}_i^{\mathrm{rec}}; b_i) \ \ (\text{all } i)
     \end{array}}
    {\Gamma \vdash \tfix\,f{:}T.\,t : T}
\end{array}
\]
\caption{Typing of terms.  In the $\tfix$ rule, $x_k$ is the designated
  \emph{guard argument}; $\tup{y}_i^{\mathrm{rec}} \subseteq \tup{y}_i$
  denotes the recursive argument positions of $\wcns{c}_i$ (those of type
  $I$); and the \emph{guard condition}
  $\mathrm{grd}_k(f; V; b)$ requires every occurrence of $f$ in $b$ to
  occur in a subterm of the form $f\,a_1 \cdots a_n$ with $a_k \in V$.
  We further require $f, \tup{x}, \tup{y}_i$ not to be shadowed in $b_i$.}
\label{fig:terms}
\end{figure}

\begin{figure}[t]
\small
\[
\begin{array}{c}
\rul{p : \varphi \in \Gamma}{\Gamma \vdash p : \varphi}
\qquad
\rul{(c \defeq \pi : \varphi) \in \Env}{\Gamma \vdash c : \varphi}
\qquad
\rul{\Gamma, p{:}\varphi \vdash \pi : \psi}
    {\Gamma \vdash \lambda p{:}\varphi.\pi : \varphi \to \psi}
\\[2.2ex]
\rul{\Gamma \vdash \pi : \varphi \to \psi \quad \Gamma \vdash \pi' : \varphi}
    {\Gamma \vdash \pi\,\pi' : \psi}
\qquad
\rul{\Gamma, x{:}\sigma \vdash \pi : \varphi}
    {\Gamma \vdash \lambda x{:}\sigma.\pi : \forall x{:}\sigma.\varphi}
\\[2.2ex]
\rul{\Gamma \vdash \pi : \forall x{:}\sigma.\varphi \quad \Gamma \vdash a : \sigma}
    {\Gamma \vdash \pi\,a : \subst{\varphi}{a}{x}}
\qquad
\rul{\Gamma, q{:}K \vdash \pi : \varphi}
    {\Gamma \vdash \Lambda q{:}K.\pi : \forall q{:}K.\varphi}
\\[2.2ex]
\rul{\Gamma \vdash \pi : \forall q{:}K.\varphi \quad \Gamma \vdash P : K}
    {\Gamma \vdash \pi\,P : \subst{\varphi}{P}{q}}
\qquad
\rul{\Gamma \vdash \pi_1 : \varphi_1 \quad \Gamma \vdash \pi_2 : \varphi_2}
    {\Gamma \vdash \langle\pi_1,\pi_2\rangle : \varphi_1 \wedge \varphi_2}
\\[2.2ex]
\rul{\Gamma \vdash \pi : \varphi_1 \wedge \varphi_2}
    {\Gamma \vdash \pi.i : \varphi_i}
\qquad
\rul{\Gamma \vdash \pi : \varphi_i}
    {\Gamma \vdash \iota_i(\pi) : \varphi_1 \vee \varphi_2}
\\[2.2ex]
\rul{\Gamma \vdash \pi : \varphi_1 \vee \varphi_2 \quad
     \Gamma, p_i{:}\varphi_i \vdash \pi_i : \psi \ (i = 1,2)}
    {\Gamma \vdash \mathsf{vcase}(\pi; p_1.\pi_1; p_2.\pi_2) : \psi}
\\[2.2ex]
\rul{\Gamma \vdash a : \sigma \quad \Gamma \vdash \pi : \subst{\varphi}{a}{x}}
    {\Gamma \vdash \langle a, \pi \rangle : \exists x{:}\sigma.\varphi}
\qquad
\rul{}{\Gamma \vdash \mathsf{tt} : \top}
\qquad
\rul{\Gamma \vdash a : \sigma}
    {\Gamma \vdash \mathsf{refl}_{\sigma}(a) : a =_{\sigma} a}
\\[2.2ex]
\rul{\Gamma \vdash \pi : \exists x{:}\sigma.\varphi \quad
     \Gamma, x{:}\sigma, p{:}\varphi \vdash \pi' : \psi \quad
     x, p \notin \fv{\psi}}
    {\Gamma \vdash \mathsf{ecase}(\pi; x\,p.\pi') : \psi}
\qquad
\rul{\Gamma \vdash \pi : \bot \quad \Gamma \vdash \varphi \sprop}
    {\Gamma \vdash \texfalso{\pi}{\varphi} : \varphi}
\\[2.2ex]
\rul{\Gamma \vdash t : \refty{x{:}\sigma}{\varphi}}
    {\Gamma \vdash \tvalid{t} : \subst{\varphi}{\tval{t}}{x}}
\qquad
\rul{\Gamma \vdash \pi : \varphi \quad \varphi \conv \varphi' \quad
     \Gamma \vdash \varphi' \sprop}
    {\Gamma \vdash \pi : \varphi'}
\\[2.2ex]
\rul{\Gamma \vdash \pi : a =_{\sigma} b \quad
     \Gamma, z{:}\sigma \vdash \psi \sprop \quad
     \Gamma \vdash \pi' : \subst{\psi}{a}{z}}
    {\Gamma \vdash \mathsf{rew}(\pi; z.\psi; \pi') : \subst{\psi}{b}{z}}
\\[2.2ex]
\rul{\begin{array}{c}
     \wcns{k} : \forall \tup{x}{:}\tup{\tau}.\,\varphi_1 \to \dots \to
       \varphi_\ell \to J(\tup{u}) \in \Env \\
     \Gamma \vdash \tup{a} : \tup{\tau} \quad
     \Gamma \vdash \pi_i : \subst{\varphi_i}{\tup{a}}{\tup{x}} \ (\text{all } i)
     \end{array}}
    {\Gamma \vdash \wcns{k}(\tup{a}; \tup{\pi}) :
       J(\subst{\tup{u}}{\tup{a}}{\tup{x}})}
\\[2.2ex]
\rul{\Gamma \vdash \pi : J(\tup{a}) \quad
     \Gamma, \tup{z}{:}\tup{\sigma} \vdash \psi \sprop \quad
     \Gamma \vdash \pi_j : \Xi_j(\tup{z}.\psi) \ \ (\text{all clauses } j)}
    {\Gamma \vdash \tind{J}(\tup{z}.\psi; \tup{\pi}; \pi) :
       \subst{\psi}{\tup{a}}{\tup{z}}}
\\[2.2ex]
\rul{\Gamma \vdash a : I \quad
     \Gamma, z{:}I \vdash \psi \sprop \quad
     \Gamma \vdash \pi_i : \Theta_i(z.\psi) \ \ (\text{all constrs } i)}
    {\Gamma \vdash \tind{I}(z.\psi; \tup{\pi}; a) : \subst{\psi}{a}{z}}
\end{array}
\]
\caption{Typing of proofs.  For a predicate clause
  $\wcns{k}_j : \forall \tup{x}{:}\tup{\tau}.\varphi_1 \to \dots \to
  \varphi_\ell \to J(\tup{u})$, the \emph{clause obligation} is
  $\Xi_j(\tup{z}.\psi) = \forall \tup{x}{:}\tup{\tau}.\,\varphi_1 \to \dots
  \to \varphi_\ell \to \subst{\psi}{\tup{v}_1}{\tup{z}} \to \dots \to
  \subst{\psi}{\tup{v}_r}{\tup{z}} \to \subst{\psi}{\tup{u}}{\tup{z}}$,
  where $J(\tup{v}_1), \dots, J(\tup{v}_r)$ are the recursive premises
  among the $\varphi_i$ (an induction hypothesis for each).  For a
  datatype constructor $\wcns{c}_i : (\tup{\sigma}_i)$ of $I$, the clause
  obligation is $\Theta_i(z.\psi) = \forall \tup{y}{:}\tup{\sigma}_i.\,
  \subst{\psi}{y_{r_1}}{z} \to \dots \to \subst{\psi}{y_{r_s}}{z} \to
  \subst{\psi}{\wcns{c}_i(\tup{y})}{z}$, where $y_{r_1}, \dots, y_{r_s}$
  are the recursive arguments.}
\label{fig:proofs}
\end{figure}

\subsection{Reduction and conversion}
\label{sec:calculus-red}

\begin{definition}[Reduction]
\label{def:source-red}
The relation $\red$ on terms is the contextual closure --- in all term,
proof, type, proposition and predicate positions --- of the following root
rules:
\begin{align*}
  (\lambda x{:}\sigma.t)\,a &\red \subst{t}{a}{x}
  & (\lambda p{:}\varphi.t)\,\pi &\red \subst{t}{\pi}{p}
  \\
  \tfix\,f{:}T.\,t &\red \subst{t}{\tfix\,f{:}T.\,t}{f}
  & \tval{\trefine{t}{\pi}} &\red t
  \\
  c &\red t \quad \text{if } (c \defeq t : \sigma) \in \Env
  & \tcast{\mathsf{refl}_\sigma(a)}{z.\tau}{t} &\red t
  \\
  \tcase{\wcns{c}_i(\tup{a})}{z.\tau}{\{\dots\}} &\red
     \subst{t_i}{\tup{a}}{\tup{y}_i}
  \\
  \tsplit{\langle \pi_1, \pi_2 \rangle}{p_1\,p_2.t} &\red
     \subst{\subst{t}{\pi_1}{p_1}}{\pi_2}{p_2}
\end{align*}
Source \emph{conversion} $\conv$ is the least congruence (on each
syntactic category, closing under all formers) containing $\red$.
\end{definition}

Reduction is untyped and unrestricted; in particular fixpoint unfolding is
not value-restricted.  We prove no metatheoretic properties of $\red$
itself (no confluence or normalisation of the source is needed anywhere in
this paper): $\conv$ enters only through the conversion typing rules, and
the semantics of \cref{sec:semantics} is invariant under it
(\cref{lem:conv-inv}).

\begin{lemma}[Rigidity of type formers]
\label{lem:rigid}
If $\sigma \conv \sigma'$ then $\sigma$ and $\sigma'$ have the same
outermost former: both are the same datatype name, or both are
$\Pi x{:}\rho.\tau$-types whose domains and codomains are respectively
convertible, or both are $\varphi \to \tau$-types with convertible
components, or both are refinement types with convertible components.  The
analogous statement holds for propositions.
\end{lemma}

\begin{proof}
No root rule of \cref{def:source-red} applies at the root of a set type or
a proposition --- the root rules rewrite terms and proofs only --- so every
$\red$-step inside a type or proposition happens strictly below its
outermost former, or below the outermost former of one of its component
types and propositions.  The claim follows by induction on the length of
the conversion path, observing that a single step preserves the outermost
former and relates components pointwise.
\end{proof}

\begin{lemma}[Substitution]
\label{lem:source-subst}
All judgements admit substitution: if
$\Gamma, x{:}\sigma, \Gamma' \vdash \mathcal{J}$ and $\Gamma \vdash a :
\sigma$, then $\Gamma, \subst{\Gamma'}{a}{x} \vdash
\subst{\mathcal{J}}{a}{x}$, and analogously for proof variables ($\Gamma
\vdash \pi : \varphi$) and predicate variables ($\Gamma \vdash P : K$).
\end{lemma}

\begin{proof}
Routine simultaneous induction on the second derivation, using that
$\conv$ is closed under substitution (each root rule is) and that
substitution into predicate atoms $q(\tup{a})[P/q] = P \cdot \tup{a}$ is
given by \emph{predicate application}: for $P = \lambda(\tup{x} {:}
\tup{\sigma}).\varphi$ we set $P \cdot \tup{a} =
\subst{\varphi}{\tup{a}}{\tup{x}}$.  Predicate application terminates:
$\varphi$ cannot contain $q$ (predicate expressions are closed with
respect to the variable being substituted after $\alpha$-renaming), so the
substitution $\subst{\cdot}{P}{q}$ is defined by structural recursion on
the phrase substituted into, with the atom case delegating to ordinary
term substitution inside $\varphi$.
\end{proof}

\subsection{The shallow fragment}
\label{sec:calculus-shallow}

\begin{definition}[Shallow fragment $\Sfrag$]
\label{def:shallow-fragment}
A proposition is \emph{shallow} if it contains no subformula of the form
$\forall q{:}K.\varphi$ and no atom $q(\tup{a})$ with $q$ a predicate
variable; a set type is shallow if all propositions occurring in it
(refinement payloads and premises) are shallow.  We write $\Sfrag$ for the
class of shallow propositions and set types.
\end{definition}

Every proposition and type built from datatypes, refinements, inductive
predicates, equality and the first-order connectives and quantifiers is
shallow; only genuine predicate quantification is excluded.  The
translation of \cref{sec:translation} is defined on $\Sfrag$; proofs
\emph{inside} the environment may use the full calculus, including
predicate quantification --- the semantics of \cref{sec:semantics} covers
it --- but a lemma whose \emph{statement} leaves $\Sfrag$ contributes no
premise axiom, a definition whose type leaves $\Sfrag$ contributes no
specification axiom, and goals must be shallow and closed.  We further
require the \emph{declarations} of $\Env$ (datatype constructor argument
types; inductive-predicate index types, clause telescopes, premises and
indices) to lie in $\Sfrag$, since the translation must emit structural
axioms for them.

\subsection{Relation to CIC}
\label{sec:calculus-cic}

$\CICm$ is designed as the image of a first-order, monomorphic slice of
CIC under analyses that are part of the CoqHammer translation pipeline or
of program extraction:

\begin{itemize}[leftmargin=2em]
\item \emph{Stratification.}  CIC's sort analysis (is this subterm a
  proof? a proposition? a set-level term?) is decidable and implemented
  (proposition detection in the translator; \cite[\S 4]{CzajkaKaliszyk2018}).
  $\CICm$ bakes its outcome into the grammar.  Proof irrelevance, which
  the translation assumes globally, is not needed in the conversion rule
  of $\CICm$ (conversion never rewrites inside a proposition's proof in a
  way visible to the semantics); it reappears as a property of the
  \emph{model} (\cref{sec:semantics}).
\item \emph{Refinement types.}  $\mathsf{sig}\ A\ P$ with its constructor
  $\mathsf{exist}$ and projections is an ordinary inductive type in CIC;
  the classification of inductive instances by propositional content
  (empty / propositional singleton / refinement / regular) identifies it,
  per instance, with the primitive $\refty{x{:}A}{P\,x}$ of $\CICm$ ---
  this is Letouzey's informative-singleton unboxing \cite[\S
  4.3]{Letouzey2004} applied at the source.  Instances of
  $\mathsf{prod}$/$\mathsf{sigT}$ with one propositional component
  classify the same way.
\item \emph{Singleton eliminations.}  CIC permits matches on
  $\Prop$-sorted scrutinees in informative positions only for empty and
  singleton propositional inductives; these appear in $\CICm$ as
  $\texfalso{}{}$, $\tsplit{}{}$ and $\tcast{}{}$ ($\bot$, $\wedge$, $=$).
\item \emph{Monomorphism.}  $\CICm$ has no type variables and no
  parameterised inductives: each datatype is declared at a ground
  instance ($\mathsf{list\ nat}$ is a datatype of its own).  Heuristic
  monomorphisation --- instantiating polymorphic definitions at the type
  instances occurring in a problem --- is an independent translation stage
  studied separately; this paper treats the post-monomorphisation
  fragment.
\item \emph{Structural fixpoints.}  The $\tfix$ former captures
  structural recursion with a designated decreasing argument, recursion on
  immediate constructor subterms, and a case-split body --- the common
  shape of $\mathsf{Fixpoint}$ definitions.  Deep pattern guards, nested
  and mutual recursion, and well-founded ($\mathsf{Acc}$-based) recursion
  are excluded; \cref{sec:limitations} discusses all four, the last being
  a genuine semantic boundary rather than a presentational one.
\item \emph{Prop-level logic.}  The primitive connectives of $\CICm$
  correspond to the inductive definitions of $\mathsf{True}$,
  $\mathsf{False}$, $\mathsf{and}$, $\mathsf{or}$, $\mathsf{ex}$ and
  $\mathsf{eq}$ in CIC; identifying them with primitives is harmless
  because their CIC induction principles are exactly the elimination
  rules of \cref{fig:proofs}.
\end{itemize}
